\section{Methodology}
\label{sec:method}

\subsection{Model and Components}
\label{sec:model}

We study \gptsmall \citep{radford2019language}, a 12-layer transformer with 12 attention heads per layer ($d_\text{model} = 768$). We evaluate all 156 components: 144 attention heads (indexed as H$_{\ell.h}$ for layer $\ell$, head $h$) and 12 MLP layers (indexed as MLP$_\ell$). We use TransformerLens for all model access and activation patching.

\subsection{Tasks and Datasets}
\label{sec:tasks}

We construct three task datasets, each with 100 clean/corrupted prompt pairs. The tasks span distinct linguistic and reasoning domains to enable meaningful cross-task selectivity evaluation.

\para{Indirect Object Identification (\ioi).} The model must predict the indirect object in sentences like ``When Anna and David went to the store, David gave a ball to \underline{\hspace{1em}}.'' The correct answer is the name that is \emph{not} the subject of the final clause (here, ``Anna'').

\para{Subject-Verb Agreement (\sva).} The model must predict the correct verb form given a subject with an intervening attractor: ``The dog above the books \underline{\hspace{1em}}'' $\rightarrow$ ``is'' (not ``are'').

\para{Greater-Than (\gt).} The model must predict a valid continuation digit for a year comparison: ``The war lasted from 1758 to 17\underline{\hspace{1em}}'' $\rightarrow$ any digit $\geq 6$.

\subsection{Corruption Strategy}
\label{sec:corruption}

Following \citet{zhang2023best}, we use symmetric token replacement (\str) rather than Gaussian noise to avoid off-distribution artifacts. For \ioi, we replace names with different names (changing the correct indirect object). For \sva, we swap singular/plural subjects and attractors (flipping the required verb number). For \gt, we change the reference year (shifting valid continuation digits). All tokens are verified as single tokens in \gptsmall's vocabulary.

\subsection{Intervention Procedures}
\label{sec:interventions}

For each component $c$ and task $t$, we compute two scores:

\para{Necessity (noising).} We run the model on clean input $x$ and replace component $c$'s activations with those from the corrupted input $\tilde{x}$. The necessity score measures how much the logit difference drops:
\begin{equation}
    \text{Nec}(c, t) = \frac{\logitdiff(x) - \logitdiff(x \mid c \leftarrow \tilde{c})}{\totaleffect}
    \label{eq:necessity}
\end{equation}
where $\totaleffect = \logitdiff(x) - \logitdiff(\tilde{x})$ is the total effect of corruption.

\para{Sufficiency (denoising).} We run the model on corrupted input $\tilde{x}$ and replace component $c$'s activations with those from the clean input $x$. The sufficiency score measures how much the logit difference recovers:
\begin{equation}
    \text{Suf}(c, t) = \frac{\logitdiff(\tilde{x} \mid c \leftarrow c^*) - \logitdiff(\tilde{x})}{\totaleffect}
    \label{eq:sufficiency}
\end{equation}
where $c^*$ denotes the clean activation.

Both scores are normalized by $\totaleffect$ to enable comparison across tasks with different effect magnitudes.

\subsection{Cross-Task Selectivity}
\label{sec:selectivity}

For each task $t$, we identify the top-$K$ components by absolute effect size ($K = 15$, approximately 10\% of all components). We then measure each component's effect on all other tasks.

\para{Selectivity score.} For a component $c$ identified as important for task $t$:
\begin{equation}
    \text{Sel}(c, t) = |\text{Effect}(c, t)| - \frac{1}{|T \setminus t|} \sum_{t' \in T \setminus t} |\text{Effect}(c, t')|
    \label{eq:selectivity}
\end{equation}
Positive values indicate the component is more specific to its target task.

\para{Cross-task leakage.} A component \emph{leaks} if its maximum absolute effect on any non-target task exceeds a threshold $\tau = 0.05$ (5\% of normalized effect). The leakage rate is the fraction of top-$K$ components that leak.

\subsection{Statistical Analysis}
\label{sec:stats}

We compare necessity and sufficiency rankings using Spearman rank correlation. We test whether the two intervention types differ in selectivity using the Mann-Whitney U test (appropriate for non-normal distributions of selectivity scores). We compute bootstrap 95\% confidence intervals for the selectivity difference (10,000 resamples) and report Cohen's $d$ for effect size.

\subsection{Hyperparameters}
\label{sec:hyperparams}

\tabref{tab:hyperparams} summarizes the experimental configuration.

\begin{table}[t]
    \centering
    \begin{tabular}{@{}llp{6cm}@{}}
        \toprule
        \textbf{Parameter} & \textbf{Value} & \textbf{Justification} \\
        \midrule
        Examples per task & 100 & Sufficient for stable patching estimates \\
        Batch size & 25 & Memory-efficient; 4 batches per task \\
        Top-$K$ & 15 & ${\sim}10\%$ of components; standard in \mi \\
        Leakage threshold $\tau$ & 0.05 & 5\% of normalized effect \\
        Random seed & 42 & Reproducibility \\
        \bottomrule
    \end{tabular}
    \caption{Experimental hyperparameters.}
    \label{tab:hyperparams}
\end{table}